%% USPSC-modelo-IFSC-DFA.tex
% ---------------------------------------------------------------
% USPSC: Modelo de Trabalho Academico (tese de doutorado, dissertacao de
% mestrado e trabalhos monograficos em geral) em conformidade com 
% ABNT NBR 14724:2011: Informacao e documentacao - Trabalhos academicos -
% Apresentacao
%----------------------------------------------------------------
%% Esta é uma customização do abntex2-modelo-trabalho-academico.tex de v-1.9.5 laurocesar 
%% para as Unidades do Campus USP de São Carlos:
%% EESC - Escola de Engenharia de São Carlos
%% IAU - Instituto de Arquitetura e Urbanismo
%% ICMC - Instituto de Ciências Matemáticas e de Computação
%% IFSC - Instituto de Física de São Carlos
%% IQSC - Instituto de Química de São Carlos
%%
%% Este trabalho utiliza a classe USPSC.cls que é mantida pela seguinte equipe:
%% 
%% Programação:
%%   - Marilza Aparecida Rodrigues Tognetti - marilza@sc.usp.br (PUSP-SC)
%%   - Ana Paula Aparecida Calabrez - aninha@sc.usp.br (PUSP-SC)
%% Normalização e Padronização:
%%   - Brianda de Oliveira Ordonho Sigolo - brianda@usp.br (IAU)
%%   - Elena Luzia Palloni Gonçalves - elena@sc.usp.br (EESC)
%%   - Eliana de Cássia Aquareli Cordeiro - eliana@iqsc.usp.br (IQSC)
%%   - Flávia Helena Cassin - cassinp@sc.usp.br (EESC)
%%   - Maria Cristina Cavarette Dziabas - mcdziaba@ifsc.usp.br (IFSC)
%%   - Regina Célia Vidal Medeiros - rcvmat@icmc.usp.br (ICMC)
%%
%% O USPSC-modelo.tex utiliza:	
%%  USPSC.cls e USPSC1.cls
%% 	USPSC-modelo-references.bib
%%	USPSC-modelo.tex
%%	USPSC-unidades.tex
%%	Um dos arquivos com dados pre-textuais abaixo, em conformidade com a Unidade de vínculo do autor:
%%				USPSC-pre-textual-EESC.tex
%%				USPSC-pre-textual-IAU.tex
%%				USPSC-pre-textual-ICMC.tex
%%				USPSC-pre-textual-IFSC.tex
%%				USPSC-pre-textual-IQSC.tex
%%				USPSC-pre-textual-OUTRO.tex
%%	USPSC-fichacatalografica.tex ou fichacatalografica.pdf
%%	folhadeaprovacao.pdf
%%	USPSC-Cap1-Introducao.tex
%%	USPSC-Cap2-Desenvolvimento.tex
%%	USPSC-Cap3-Citacoes.tex
%%	USPSC-Cap4-referencias.tex
%%	USPSC-Cap5-Conclusao.tex
%%	USPSC-Apendices.tex
%%	USPSC-Anexos.tex
%%	USPSC-AcentuacaoLaTeX.tex
%%	USPSC-LetrasGregas.tex
%%	USPSC-SimbolosUteis.tex

%----------------------------------------------------------------
%% Sobre a classe abntex2.cls:
%% abntex2.cls, v-1.9.5 laurocesar
%% Copyright 2012-2015 by abnTeX2 group at https://www.abntex.net.br/ 
%%
%----------------------------------------------------------------

\documentclass[
% -- opções da classe memoir --
12pt,		% tamanho da fonte
openright,	% capítulos começam em pág ímpar (insere página vazia caso preciso)
twoside,  % para impressão em anverso (frente) e verso. Oposto a oneside - Nota: utilizar \imprimirfolhaderosto*
%oneside, % para impressão em páginas separadas (somente anverso) -  Nota: utilizar \imprimirfolhaderosto
% inclua uma % antes do comando twoside e exclua a % antes do oneside 
a4paper,			% tamanho do papel. 
% -- opções da classe abntex2 --
chapter=TITLE,		% títulos de capítulos convertidos em letras maiúsculas
% -- opções do pacote babel --
english,			% idioma adicional para hifenização
french,				% idioma adicional para hifenização
spanish,			% idioma adicional para hifenização
brazil				% o último idioma é o principal do documento
% {USPSC} configura o cabeçalho contendo apenas o número da página
]{USPSC}
%]{USPSC1}
% Inclua % antes de ]{USPSC} e retire a % antes de %]{USPSC1}
% para utilizar o cabeçalho diferenciado para as páginas pares e ímpares como indicado abaixo:
%- páginas ímpares: cabeçalho com seções ou subseções e o número da página
%- páginas pares: cabeçalho com o número da página e o título do capítulo 
% ---

% ---
% Pacotes básicos - Fundamentais 
% ---
\usepackage[T1]{fontenc}		% Seleção de códigos de fonte.
\usepackage[utf8]{inputenc}		% Codificação do documento (conversão automática dos acentos)
\usepackage{lmodern}			% Usa a fonte Latin Modern
% Para utilizar a fonte Times New Roman, inclua uma % no início do comando acima  "\usepackage{lmodern}"
% Abaixo, tire a % antes do comando  \usepackage{times}
%\usepackage{times}		    	% Usa a fonte Times New Roman	
% Lembre-se de alterar a fonte no comando que imprime o preâmbulo no arquivo da Classe USPSC.cls				
\usepackage{lastpage}			% Usado pela Ficha catalográfica
\usepackage{indentfirst}		% Indenta o primeiro parágrafo de cada seção.
\usepackage{color}				% Controle das cores
\usepackage{graphicx}			% Inclusão de gráficos
\usepackage{float} 				% Fixa tabelas e figuras no local exato
\usepackage{chemfig,chemmacros} % Para escrever reações químicas
\usepackage{microtype} 			% para melhorias de justificação
\usepackage{pdfpages}
\usepackage{makeidx}            % para gerar índice remissimo
\usepackage{listings}
\usepackage{color}
\usepackage{graphicx}		
\usepackage{amsmath}	
\usepackage{microtype}
\usepackage{bookmark}


\lstset{frame=tb,
	language=C++,
	aboveskip=3mm,
	belowskip=3mm,
	showstringspaces=false,
	columns=flexible,
	basicstyle={\small\ttfamily},
	numbers=none,
	numberstyle=\tiny\color{gray},
	keywordstyle=\color{blue},
	commentstyle=\color{green},
	stringstyle=\color{red},
	breaklines=true,
	breakatwhitespace=true,
	tabsize=3
}
% ---

% ---
% Pacotes de citações
% Citações padrão ABNT
% ---
% Sistemas de chamada: autor-data ou numérico.
% Sistema autor-data
%\usepackage[alf,abnt-emphasize=bf, abnt-thesis-year=both, abnt-repeated-author-omit=yes, abnt-last-names=abnt, abnt-etal-cite,abnt-etal-list=3, abnt-etal-text=default, abnt-and-type=e, abnt-doi=doi, abnt-url-package=none, abnt-verbatim-entry=no]{abntex2cite}

% Para o IQSC, que indica todos os autores nas referências, incluir % no início do comando acima e retirar a % do comando abaixo 

%\usepackage[alf,abnt-emphasize=bf, abnt-thesis-year=both, abnt-repeated-author-omit=yes, abnt-last-names=abnt, abnt-etal-cite,abnt-etal-list=0, abnt-etal-text=default, abnt-and-type=e]{abntex2cite}

% Sistema Numérico
% Para citações numéricas, sistema adotado pelo IFSC, incluir % no início do comando acima e retirar a % do comando abaixo 
\usepackage[num,overcite,abnt-emphasize=bf, abnt-thesis-year=both, abnt-repeated-author-omit=yes, abnt-last-names=abnt, abnt-etal-cite,abnt-etal-list=0, abnt-etal-text=default, abnt-and-type=e]{abntex2cite}

% Complementarmente, verifique as instruções abaixo sobre os Pacotes de Nota de rodapé
% ---
% Pacotes de Nota de rodapé
% Configurações de nota de rodapé

% O presente modelo adota o formato numérico para as notas de rodapés quando utiliza o sistema de chamada autor-data para citações e referências. Para utilizar o sistema de chamada numérico para citações e referências, habilitar um dos comandos abaixo.
% Há diversa opções para nota de rodapé no Sistema Numérico.  Para o IFSC, habilitade o comando abaixo.

\renewcommand{\thefootnote}{\fnsymbol{footnote}}  %Comando para inserção de símbolos em nota de rodapé

% Outras opções para nota de rodapé no Sistema Numérico:
%\renewcommand{\thefootnote}{\alph{footnote}}      %Comando para inserção de letras minúscula em nota de rodapé
%\renewcommand{\thefootnote}{\Alph{footnote}}      %Comando para inserção de letras maiúscula em nota de rodapé
%\renewcommand{\thefootnote}{\roman{footnote}}     %Comando para inserção de números romanos minúsculos  em nota de rodapé
%\renewcommand{\thefootnote}{\Roman{footnote}}     %Comando para inserção de números romanos minúsculos  em nota de rodapé

\renewcommand{\footnotesize}{\small} %Comando para diminuir a fonte das notas de rodapé

 % ---
 % Pacote para agrupar a citação numérica consecutiva
 % Quando for adotado o Sistema Numérico, a exemplo do IFSC, habilite 
 % o pacote cite abaixo retirando a porcentagem antes do comando abaixo
 %\usepackage[superscript]{cite}	

% ---
% Pacotes adicionais, usados apenas no âmbito do Modelo Canônico do abnteX2
% ---
\usepackage{lipsum}				% para geração de dummy text
% ---

\usepackage{caption}
\captionsetup{format= hang,labelsep=endash,labelfont=bf} 

% pacotes de tabelas
\usepackage{multicol}	% Suporte a mesclagens em colunas
\usepackage{multirow}	% Suporte a mesclagens em linhas
\usepackage{longtable}	% Tabelas com várias páginas
\usepackage{threeparttablex}    % notas no longtable
\usepackage{array}
\usepackage{subcaption}
%---
% Configurações para o pacote chemfig
%\chemsetup[chemformula]{format=\sffamily}
\renewcommand*\printatom[1]{\ensuremath{\mathsf{#1}}}
%\setatomsep{2em}
%\setdoublesep{.6ex}
%\setbondstyle{semithick}
%---
% Configurando o ambiente para utilizar os recursos de frases pre-prontas do mhchem
\newenvironment{rslist}%
{%
	\begin{labeling}% environment from KOMA-script
		{\rsnumber{R39/23/24/25}}% R39/23/24/25 is longest label
	}{%
\end{labeling}%
}%
% Definição de comando para utilizar os recursos de frases pre-prontas do mhchem
\newcommand{\rs}[2][]{\item[{\rsnumber[#1]{#2}}] \rsphrase{bb}}
% ---

% ---
% DADOS INICIAIS - Define sigla com título, área de concentração e opção do Programa 
% Consulte a tabela referente aos Programas, áreas e opções de sua unidade contante do
% arquivo USPSC-Siglas estabelecidas para os Programas de Pós-Graduação por Unidade.xlsx 
% ou nos APÊNDICES A-G
\siglaunidade{IFSC} 
\programa{DFAFC}
% Os demais dados deverão ser fornecidos no arquivo USPSC-pre-textual-UUUU ou USPSC-TCC-pre-textual-UUUU, onde UUUU é a sigla da Unidade. 
% Exemplo:USPSC-pre-textual-IFSC.tex
% ---
% Configurações de aparência do PDF final
% alterando o aspecto da cor azul
\definecolor{blue}{RGB}{41,5,195}

% informações do PDF
\makeatletter
\hypersetup{
	%pagebackref=true,
	pdftitle={\@title}, 
	pdfauthor={\@author},
	pdfsubject={\imprimirpreambulo},
	pdfcreator={LaTeX with abnTeX2},
	pdfkeywords={abnt}{latex}{abntex}{USPSC}{trabalho acadêmico}, 
	colorlinks=true,       		% false: boxed links; true: colored links
	linkcolor=blue,          	% color of internal links
	citecolor=blue,        		% color of links to bibliography
	filecolor=magenta,      		% color of file links
	urlcolor=blue,
	bookmarksdepth=4
}
\makeatother
% --- 
% --- 
% Espaçamentos entre linhas e parágrafos 
% --- 

% O tamanho do parágrafo é dado por:
\setlength{\parindent}{1.3cm}

% Controle do espaçamento entre um parágrafo e outro:
\setlength{\parskip}{0.2cm}  % tente também \onelineskip

% ---
% compila o sumário e índice
\makeindex
% ---

% ----
% Início do documento
% ----
\begin{document}

% Seleciona o idioma do documento (conforme pacotes do babel)
\selectlanguage{brazil}
% Se o idioma do texto for inglês, inclua uma % antes do 
%      comando \selectlanguage{brazil} e 
%      retire a % antes do comando abaixo
%\selectlanguage{english}

% Retira espaço extra obsoleto entre as frases.
\frenchspacing 

% --- Formatação dos Títulos
\renewcommand{\ABNTEXchapterfontsize}{\fontsize{12}{12}\bfseries}
\renewcommand{\ABNTEXsectionfontsize}{\fontsize{12}{12}\bfseries}
\renewcommand{\ABNTEXsubsectionfontsize}{\fontsize{12}{12}\normalfont}
\renewcommand{\ABNTEXsubsubsectionfontsize}{\fontsize{12}{12}\normalfont}
\renewcommand{\ABNTEXsubsubsubsectionfontsize}{\fontsize{12}{12}\normalfont}


% ----------------------------------------------------------
% ELEMENTOS PRÉ-TEXTUAIS
% ----------------------------------------------------------
% ---
% Capa
% ---
\imprimircapa
% ---
% Folha de rosto
% (o * indica impressão em anverso (frente) e verso )
% ---
\imprimirfolhaderosto*
%\imprimirfolhaderosto
% ---
% ---
% Inserir a ficha catalográfica em pdf
% ---
% A biblioteca da sua Unidade lhe fornecerá um PDF com a ficha
% catalográfica definitiva. 
% Quando estiver com o documento, salve-o como PDF no diretório
% do seu projeto como fichacatalografica.pdf e iclua o arquivo
% utilizando o comando abaixo:
%\includepdf{USPSC-Tese-pre-textual/fichacatalografica.pdf} 
% Se você optar por elaborar a ficha catalográfica, deverá 
% incluir uma % antes da linha % antes
% do comando % ---
% Inserir a ficha bibliografica
% ---
% Isto é um exemplo de Ficha Catalográfica, ou ``Dados internacionais de
% catalogação-na-publicação''. Você pode utilizar este modelo como referência. 
% Porém, provavelmente a biblioteca da sua universidade lhe fornecerá um PDF
% com a ficha catalográfica definitiva após a defesa do trabalho. Quando estiver
% com o documento, salve-o como PDF no diretório do seu projeto e substitua todo
% o conteúdo de implementação deste arquivo pelo comando abaixo:
%
\begin{fichacatalografica}
	\hspace{-1.4cm}
	\imprimirnotaautorizacao \\ \\
	%\sffamily
	\vspace*{\fill}					% Posição vertical
	\begin{center}					% Minipage Centralizado
		\imprimirnotabib \\
\begin{table}[htb]
	\scriptsize
	\centering	
	\begin{tabular}{|p{0.9cm} p{8.7cm}|}
		\hline
	      & \\
		  &	  \imprimirautorficha     \\
		
		 \imprimircutter & 
							\hspace{0.4cm}\imprimirtitulo~  / ~\imprimirautor~ ;  ~\imprimirorientadorcorpoficha. -- 	\imprimirlocal, \imprimirdata.   \\
		
		  &  % Para incluir nota referente à versão corrigida no corpo da ficha,
			  % incluir % no início da linha acima e tirar a % do início da linha abaixo
			  %	\hspace{0.4cm} \imprimirtitulo~  / ~\imprimirautor~ ; ~\imprimirorientadorcorpoficha~- ~\imprimirnotafolharosto. -- \imprimirlocal, \imprimirdata.  \\
		
			\hspace{0.4cm}\pageref{LastPage} p. : il. (algumas color.) ; 30 cm.\\ 
		  & \\
		  & 
		    \hspace{0.4cm}\imprimirnotaficha ~--~ 
						  \imprimirunidademin, 
						  \imprimiruniversidademin, 
		                  \imprimirdata. \\ 
		  & \\                 
		   % Para incluir nota referente à versão corrigida em notas,
		    % incluir uma % no início da linha acima e	
		    % tirar a % do início da linha abaixo
		    % & \hspace{0.4cm}\imprimirnotafolharosto \\ 
		  & \\ 
		  & \hspace{0.4cm}1. LaTeX. 2. abnTeX. 3. Classe USPSC. 4. Editoração de texto. 5. Normalização da documentação. 6. Tese. 7. Dissertação. 8. Documentos (elaboração). 9. Documentos eletrônicos. I. \imprimirorientadorficha. 
		   II. Título. \\
	
		     %Se houver co-orientador, inclua % antes da linha (antes de II. Título.) 
		     %          e tire a % antes do comando abaixo 
		     %III. Título. \\   
		  \hline
	\end{tabular}
\end{table}
	\end{center}
\end{fichacatalografica}
% ---

 
% e retirar o % do comando abaixo
%% ---
% Inserir a ficha bibliografica
% ---
% Isto é um exemplo de Ficha Catalográfica, ou ``Dados internacionais de
% catalogação-na-publicação''. Você pode utilizar este modelo como referência. 
% Porém, provavelmente a biblioteca da sua universidade lhe fornecerá um PDF
% com a ficha catalográfica definitiva após a defesa do trabalho. Quando estiver
% com o documento, salve-o como PDF no diretório do seu projeto e substitua todo
% o conteúdo de implementação deste arquivo pelo comando abaixo:
%
\begin{fichacatalografica}
	\hspace{-1.4cm}
	\imprimirnotaautorizacao \\ \\
	%\sffamily
	\vspace*{\fill}					% Posição vertical
	\begin{center}					% Minipage Centralizado
		\imprimirnotabib \\
\begin{table}[htb]
	\scriptsize
	\centering	
	\begin{tabular}{|p{0.9cm} p{8.7cm}|}
		\hline
	      & \\
		  &	  \imprimirautorficha     \\
		
		 \imprimircutter & 
							\hspace{0.4cm}\imprimirtitulo~  / ~\imprimirautor~ ;  ~\imprimirorientadorcorpoficha. -- 	\imprimirlocal, \imprimirdata.   \\
		
		  &  % Para incluir nota referente à versão corrigida no corpo da ficha,
			  % incluir % no início da linha acima e tirar a % do início da linha abaixo
			  %	\hspace{0.4cm} \imprimirtitulo~  / ~\imprimirautor~ ; ~\imprimirorientadorcorpoficha~- ~\imprimirnotafolharosto. -- \imprimirlocal, \imprimirdata.  \\
		
			\hspace{0.4cm}\pageref{LastPage} p. : il. (algumas color.) ; 30 cm.\\ 
		  & \\
		  & 
		    \hspace{0.4cm}\imprimirnotaficha ~--~ 
						  \imprimirunidademin, 
						  \imprimiruniversidademin, 
		                  \imprimirdata. \\ 
		  & \\                 
		   % Para incluir nota referente à versão corrigida em notas,
		    % incluir uma % no início da linha acima e	
		    % tirar a % do início da linha abaixo
		    % & \hspace{0.4cm}\imprimirnotafolharosto \\ 
		  & \\ 
		  & \hspace{0.4cm}1. LaTeX. 2. abnTeX. 3. Classe USPSC. 4. Editoração de texto. 5. Normalização da documentação. 6. Tese. 7. Dissertação. 8. Documentos (elaboração). 9. Documentos eletrônicos. I. \imprimirorientadorficha. 
		   II. Título. \\
	
		     %Se houver co-orientador, inclua % antes da linha (antes de II. Título.) 
		     %          e tire a % antes do comando abaixo 
		     %III. Título. \\   
		  \hline
	\end{tabular}
\end{table}
	\end{center}
\end{fichacatalografica}
% ---


% As informações que compõem a ficha catalográfica estão 
% definidos no arquivo USPSC-pre-textual-UUUU.tex
% ---


% ---
% ---
% Inserir errata
% ---

%\begin{errata}
%	\OnehalfSpacing 			
%	A errata é um elemento opcional, que consiste de uma lista de erros da obra, precedidos pelas folhas e linhas onde eles ocorrem e seguidos pelas correções correspondentes. Deve ser inserida logo após a folha de rosto e conter a referência do trabalho para facilitar sua identificação, conforme a ABNT NBR 14724 \cite{nbr14724}.
	
%	Modelo de Errata:
		
%	\begin{flushleft} 
%			\setlength{\absparsep}{0pt} % ajusta o espaçamento da referência	
%			\SingleSpacing 
%			\imprimirautorabr~ ~\textbf{\imprimirtitulo}.	\imprimirdata. \pageref{LastPage}p. 
			%Substitua p. por f. quando utilizar oneside em \documentclass
			%\pageref{LastPage}f.
%			\imprimirtipotrabalho~-~\imprimirinstituicao, \imprimirlocal, \imprimirdata. 
% 	\end{flushleft}
%\vspace{\onelineskip}
%\OnehalfSpacing 
%\center
%\textbf{ERRATA}
%\vspace{\onelineskip}
%\OnehalfSpacing 
%\begin{table}[htb]
%	\center%
%	\footnotesize
%	\begin{tabular}{p{1.4cm} p{1cm} p{3cm} p{3cm} }
%		\hline
%		\textbf{Folha} & \textbf{Linha}  & \textbf{Onde se lê}  & \textbf{Leia-se}  \\
%			\hline
%			1 & 10 & auto-conclavo & autoconclavo\\
%		\hline
%	\end{tabular}
%\end{table}


%x\end{errata}
% ---

% ---
% Inserir folha de aprovação
% ---

% A Folha de aprovação é um elemento obrigatório da NBR 4724/2011 (seção 4.2.1.3). 
% Após a defesa/aprovação do trabalho, gere o arquivo folhadeaprovacao.pdf da página assinada pela banca 
% e iclua o arquivo utilizando o comando abaixo:
%\includepdf{ficha_corrigida.pdf}
%\includepdf{folhadeaprovacao.pdf}
% Alternativa para a Folha de Aprovação:
% Se for a sua opção elaborar uma folha de aprovação, insira uma % antes do comando acima que inclui o arquivo folhadeaprovacao.pdf,
% tire o % do comando abaixo e altere o arquivo folhadeaprovacao.tex conforme suas necessidades
%% Alternativa para a Folha de Aprova��o 
% Se esta for a sua op��o, exclua inclus�o feita acima do arquivo folhadeaprovacao.pdf
%
\begin{folhadeaprovacao}
  \begin{center}
       {\ABNTEXchapterfont\bfseries\large\imprimirautor}
	 \vspace*{2cm}
   
    \begin{center}
      \ABNTEXchapterfont\bfseries\Large\imprimirtitulo
    \end{center}
		\vspace*{2cm}
		\hspace{.45\textwidth}
    \begin{minipage}{.5\textwidth}
        \imprimirpreambulo
    \end{minipage}
		\vspace*{2cm}
    %\vspace*{\fill}
	\end{center}

  \begin{center}
	  {\ABNTEXchapterfont\bfseries\large\ {Data de defesa: 02 de outubro de 2015} \\}
		\vspace*{\fill}
	  {\ABNTEXchapterfont\bfseries\large\ {Comiss\~ao Julgadora:} \\}
		%Trabalho aprovado. \imprimirlocal, 2 de outubro de 2015:
		
		%\assinatura{\textbf{\imprimirorientador} \\ Orientador} 
		\assinatura{\textbf{\imprimirorientador} \\}
    % Se for ORIENTADOR, inclua % no in�cio do comando abaixo e tire a % do pr�ximo comando 
		\renewcommand{\orientadorname}{Orientadora}
		%\renewcommand{\orientadorname}{Orientador}

		\imprimirorientadorRotulo
		\par
		\assinatura{\textbf{Professor} \\ Convidado1}
	
		\assinatura{\textbf{Professor} \\ Convidado2}
		%\assinatura{\textbf{Professor} \\ Convidado3}
		%\assinatura{\textbf{Professor} \\ Convidado4}
		%\begin{center}
		\vspace*{\fill}
   	{\ABNTEXchapterfont\bfseries\large\imprimirlocal}
    \par
    {\ABNTEXchapterfont\bfseries\large\imprimirdata}
\end{center}
\end{folhadeaprovacao}
% ---

%\includepdf{PaginaEmBranco.pdf}

% ---
% Dedicatória
% ---
%%% USPSC-Dedicatória.tex
% ---
% Dedicatória
% ---
\begin{dedicatoria}
   \vspace*{\fill}
   \centering
   \noindent
   \textit{ Este trabalho é dedicado aos meus pais Claudia e Geraldo.} \vspace*{\fill}
\end{dedicatoria}
% ---
% ---

% ---
% Agradecimentos
% ---
%%% Agradecimentos.tex
% ---
% Agradecimentos
% ---=====
\begin{agradecimentos}
	agradecimentos		
\end{agradecimentos}
% ---

% ---

% ---
% Epígrafe
% ---
%\begin{epigrafe}
%    \vspace*{\fill}
%	\begin{flushright}
%		\textit{"It gets easier.\\
%			Everyday, it gets a little easier.\\
%			But you gotta do it everyday.\\
%			That's the hard part.\\
%			But it does get easier." \\
%			-Jogging Baboon from BoJack Horseman
%			}
%	\end{flushright}
%\end{epigrafe}
% ---

% A T E N Ç Ã O
% Se o idioma do texto for em inglês, o abstract deve preceder o resumo
% resumo em português
%
% Resumo
% ---
\newpage\null\thispagestyle{empty}\newpage
%% Resumo.tex
% ---
% Resumo
% ---
\setlength{\absparsep}{18pt} % ajusta o espaçamento dos parágrafos do resumo		
\begin{resumo}
	\begin{flushleft} 
			\setlength{\absparsep}{0pt} % ajusta o espaçamento da referência	
			\SingleSpacing 
			\imprimirautorabr~ ~\textbf{Titulo}.	\imprimirdata. \pageref{LastPage}p. 
			%Substitua p. por f. quando utilizar oneside em \documentclass
			%\pageref{LastPage}f.
			%\imprimirtipotrabalho~-~\imprimirinstituicao, \imprimirlocal, \imprimirdata. 
			Qualificação (Doutorado em Ciências) -~\imprimirinstituicao, \imprimirlocal, \imprimirdata.
 	\end{flushleft}
\OnehalfSpacing

	texto

 \textbf{Palavras-chave}: Area importante. Grande Area. Pequena Area.

\end{resumo}

\clearpage
% ---

% Abstract
% ---
%%% Abstract.tex
% ---
% Abstract
% ---
\autor{Silva, M. J.}
\begin{resumo}[Abstract]
 \begin{otherlanguage*}{english}
	\begin{flushleft} 
		\setlength{\absparsep}{0pt} % ajusta o espaçamento dos parágrafos do resumo		
 		\SingleSpacing 
		\imprimirautorabr~ ~\textbf{Title}	\imprimirdata.  \pageref{LastPage}p. 
		%Substitua p. por f. quando utilizar oneside em \documentclass
		%\pageref{LastPage}f.
		\imprimirtipotrabalho~-~\imprimirinstituicao, \imprimirlocal, 	\imprimirdata. 
 	\end{flushleft}
	\OnehalfSpacing 
	text
   \vspace{\onelineskip}
 
   \noindent 
   \textbf{Keywords}: Important field. Great field. Small field.

 \end{otherlanguage*}
\end{resumo}

% ---

% ---
% inserir lista de figurass
% ---
%\pdfbookmark[0]{\listfigurename}{lof}
%\listoffigures*
%\cleardoublepage
% ---

% ---
% inserir lista de tabelas
% ---
%\pdfbookmark[0]{\listtablename}{lot}
%\listoftables*
%\cleardoublepage
% ---

% ---
% inserir lista de quadros
% ---
%\pdfbookmark[0]{\listofquadroname}{loq}
%\listofquadro*
%\cleardoublepage
% ---

% ---
% inserir lista de abreviaturas e siglas
% ---
%\begin{siglas}
%	\item[ANP] Agência Nacional do Petróleo, Gás Natural e Biocombustíveis
%\end{siglas}
% ---

% ---
% inserir lista de símbolos
% ---
%\begin{simbolos}
%	\item[$ m $] Massa
%\end{simbolos}
% ---
% ---
% inserir o sumario
% ---
\pdfbookmark[0]{\contentsname}{toc}
\tableofcontents*
\cleardoublepage
% ---
% ----------------------------------------------------------
% ELEMENTOS TEXTUAIS
% ----------------------------------------------------------
\textual
% Os capítulos são inseridos como arquivos externos 


\chapter{Introdução}
\label{chapter:introducao}
O openfoam \cite{Weller1998}.



\chapter{Métodos}
\label{chapter:metodos}
texto



\chapter{Resultados}
\label{chapter:resultados}
texto


\chapter{Cronograma de desenvolvimento do projeto}
\label{chapter:cronograma}

Uma vez que os resultados com a ... foram satisfatórios, na segunda metade do ... focaremos em .... O restante deste projeto será executado seguindo as seguintes fases:
%indicar mudança para modelos biológicos
\begin{enumerate}
    \item \textbf{Fase X:} Buscaremos referências na literatura dos modelos hemodinâmicos mais recentes para o estudo de variações morfológicas ou patologias no sistema cardiovascular.
    
    \item \textbf{Fase Y:} Esta fase consiste na aplicação e adaptação do fluxo de trabalho apresentado na Figura \ref{Fig:04} e discutido no Capítulo \ref{chapter:metodos}.
    
\end{enumerate}

Por se tratar de um projeto de ..., este cronograma definindo as atividades a serem feitas pós ... é subdividido em 4 semestres. As fases propostas para o plano de trabalho estão organizadas na Tabela \ref{tab:01}.

\begin{table}[h]
\centering
\caption{Cronograma das atividades.}
\begin{tabular}{|c|c|c|c|c|}
\hline
\multirow{2}{*}{{Etapas}} & \multicolumn{4}{c|}{{Semestres}}
\\ \cline{2-5}
    & 5 & 6 & 7 & 8 
\\ \hline
    {1. Fase X} & x & x & & 
\\ \hline
    {2. Fase Y} & & x & x & 
\\ \hline
\end{tabular}
\label{tab:01}
\legend{\textbf{Fonte:} Elaborada pelo autor.}
\end{table}


\chapter{Considerações Finais}
\label{chapter:conclusao}
texto


% ----------------------------------------------------------
% ELEMENTOS PÓS-TEXTUAIS
% ----------------------------------------------------------
\postextual
% ----------------------------------------------------------

% -----------------------------------------------------------
% Referências bibliográficas
% ----------------------------------------------------------
\bibliography{bibliography}


% ----------------------------------------------------------
% Glossário
% ----------------------------------------------------------
%
% Consulte o manual da classe abntex2 para orientações sobre o glossário.
%
%\glossary

% ----------------------------------------------------------
% Apêndices
% ----------------------------------------------------------


\begin{apendicesenv}
	% Imprime uma página indicando o início dos apêndices
	\partapendices
	
	\chapter{Apendice}
	\label{apendice:A}
	Texto 

	
	
\end{apendicesenv}
% ----------------------------------------------------------
% Anexos
% ----------------------------------------------------------
%\include{USPSC-Anexos}

%---------------------------------------------------------------------
% INDICE REMISSIVO
%--------------------------------------------------------------------
\include{USPSC-IndicesRemissivos}

%---------------------------------------------------------------------

\end{document}
